\documentclass[10pt,twocolumn,letterpaper]{article}

% --- Packages ---
\usepackage[utf8]{inputenc}
\usepackage[T1]{fontenc}
\usepackage{times}
\usepackage{graphicx}
\usepackage{amsmath,amssymb}
\usepackage{booktabs}
\usepackage[margin=0.75in]{geometry}
\usepackage{algorithm}
\usepackage{algpseudocode}
\usepackage{xcolor}
\usepackage{hyperref}
\usepackage{multirow}

% --- Page setup ---
\setlength{\columnsep}{0.25in}
\pagestyle{plain}
\graphicspath{{./figures/}}

% --- Custom commands ---
\newcommand{\RAMP}{\textsc{Ramp}}
\newcommand{\RRN}{\textsc{Rrn}}
\newcommand{\RRE}{RRE}
\newcommand{\CFR}{\text{CFR}}
\newcommand{\PLE}{\text{PLE}}

% --- Title ---
\title{\Large\bf Reference Reliability Estimation for\\Visual Prompting Pipelines}

\author{Anonymous Authors\\[2pt]
\texttt{\{anonymous\}@anonymous.edu}}

\date{}

\begin{document}
\maketitle

% ======================================================================
\begin{abstract}
Visual prompting pipelines---a frozen matcher feeding point prompts to a
frozen segmenter---have become a standard approach for few-shot
segmentation, yet they lack a fundamental component: \emph{reference
reliability estimation}. Swapping the reference image can change the Dice
score from 0.82 to 0.02, and the pipeline offers no signal about which
outcome occurred. We formalize this as the \textbf{Reference Reliability
Estimation} (RRE) problem: given a query and a set of candidate references,
predict the probability that each reference will produce a successful
segmentation. We show that two failure
modes---reference-induced instability and catastrophic failure---exhibit a
reproducible empirical association across four segmenters and two matchers
(Pearson $r=-0.49$ for SAM-B, $p<0.001$), and we trace this
association to the geometry of the point-prompt bottleneck through
three geometric hypotheses linking matching-map entropy to prompt dispersion
to catastrophic risk. Building on this analysis, we propose a
two-tier solution: \RAMP{} (Reference-Adaptive Multi-Prompt Fusion), a
dual-model selective predictor using a lightweight supervised risk estimator
that reduces the catastrophic failure rate by 37.5--84.4\% at 80\% coverage
on MoNuSAC, and \RRN{} (Reference Reliability Network), a lightweight learned
estimator that maps five pipeline-level features to a reliability
score, achieving AUROC = 0.612 for per-reference success prediction and
Dice@Top-3 = 0.525 on the reference selection benchmark. We establish the
first reference selection benchmark, evaluating five strategies (random,
similarity, entropy, \RRN{}, and oracle) across medical and natural image
datasets. A simple frozen top-$K$ pipeline (DINOv2-L + SAM-B, Dice = 0.446)
outperforms five recent methods---GF-SAM (0.388), MAUP
(0.279), CycleSAM (0.412), H-SAM (0.405), and SAM-Med2D (0.438)---under
a shared reference-resampling protocol, demonstrating that reference reliability
is a dimension orthogonal to prompt sophistication. The
\RAMP{} selective predictor uses coverage-based ranking that does not require
a fixed probability threshold; \RRN{} adds a lightweight learned component that narrows the
gap between similarity-based selection and the oracle upper bound.
\end{abstract}

% ======================================================================
\section{Introduction}

A user performing few-shot segmentation faces a decision that the
literature has largely ignored: which reference image should they choose?
The answer matters more than most practitioners realize. On MoNuSAC, a
query image can obtain Dice~=~0.82 with one reference and Dice~=~0.02 with
another. The pipeline provides no mechanism to distinguish these outcomes.
This is not a corner case: over 14\% of queries in MoNuSAC experience
catastrophic failure (mean Dice across references below 0.1) from at least
one reference. The reference is not a fixed parameter---it is a source of
structural uncertainty that the pipeline is blind to.

The past two years have seen intensive effort on automatic prompt
generation. GF-SAM~\cite{zhang2024gfsam} (NeurIPS 2024 Spotlight) uses
graph-based point clustering; MAUP~\cite{zhu2025maup} (MICCAI 2025) performs
K-means on confidence maps; SynPo~\cite{liu2025synpo} (MICCAI 2025) fuses
DINOv2 and SAM confidence; Matcher~\cite{liu2024matcher} (ICLR 2024) uses
bidirectional patch matching; PerSAM~\cite{zhang2024persam} (ICLR 2024)
injects target-guided attention; CycleSAM~\cite{hou2024cyclesam} (CVAI 2024)
enforces spatial cycle-consistency; H-SAM~\cite{wu2025hisam} (TPAMI 2025)
extends SAM hierarchically; and SAM-Med2D~\cite{cheng2023sammed2d} adapts
SAM for medical images. All optimize \emph{how} to construct prompts from a
given reference. None addresses \emph{whether} the reference is reliable.

This paper identifies a missing component in the visual prompting paradigm:
\textbf{Reference Reliability Estimation (RRE)}. RRE asks a simple but
consequential question: given a query and a candidate reference, what is the
probability that the reference will produce a successful segmentation? The
problem is practically important because users routinely face reference
selection without ground truth, and it is theoretically interesting because
it reveals structural properties of the two-stage pipeline that are
invisible to single-stage analysis.

\begin{figure}[t]
\centering
\includegraphics[width=\columnwidth]{fig1_anticorrelation_scatter}
\caption{The reference failure landscape. Instability (best-worst Dice gap)
  vs.\ catastrophic failure rate across four segmenters on MoNuSAC. Each
  point is a query. The structural anti-correlation---low-instability
  queries tend toward catastrophic failure, high-instability queries tend
  away from it---is a universal property of the two-stage pipeline,
  explained by the entropy of the matching map (Section~\ref{sec:theory}).}
\label{fig:anticorrelation}
\end{figure}

The central insight of this paper is that reference-induced failure has
\emph{structure}. Instability (large Dice variance across references) and
catastrophic failure (mean Dice~$<$~0.1) are not independent: they are
anti-correlated. A query that produces highly variable results across
references rarely fails catastrophically; a query that fails catastrophically
shows almost no variance across references. This pattern holds across four
segmenters---SAM-B, SAM-L, SAM2-Tiny, and SAM2-Small---and
two matchers (Fig.~\ref{fig:anticorrelation}). It is a property of the
pipeline architecture, not of any particular model.

We trace why this association exists through the geometry of the
point-prompt bottleneck (Section~\ref{sec:theory}). The key quantity is the
entropy of the matching map produced by the matcher. When entropy is low,
the matching map has sharp peaks; point coordinates are tightly clustered
around the target, and changing the reference shifts the peaks slightly
while keeping points on-target---producing high instability but low
catastrophic risk. When entropy is high, the map is diffuse; points are
scattered, and reference changes have negligible effect on already-random
coordinates---producing low instability but high catastrophic risk. We
formalize this mechanism through three geometric hypotheses linking matching
entropy to prompt dispersion, prompt dispersion to point localization error,
and point localization error to catastrophic risk---each accompanied by
empirical verification and validity controls to rule out metric coupling
artifacts.

We address the RRE problem with a two-tier solution. \textbf{\RAMP{}
(Reference-Adaptive Multi-Prompt Fusion)} is a dual-model selective
predictor that uses a lightweight supervised risk estimator (gradient-boosted
trees) with four test-time features to estimate
instability and catastrophic risk separately. It weights references by
predicted reliability, adapts prompt configuration to the risk level, and
fuses predictions via risk-stratified consensus. The result is a
selective predictor that reduces the catastrophic failure
rate by 37.5--84.4\% at 80\% coverage on MoNuSAC, using coverage-based
ranking that does not require a fixed probability threshold. \textbf{\RRN{} (Reference Reliability Network)} is a lightweight learned estimator
that takes five pipeline-level features---matching entropy, feature
similarity, prompt dispersion, mask consistency, and prediction
uncertainty---and outputs a reliability score $P(\text{success} \mid q, r)$. \RRN{}
is trained on the existing multi-reference data with binary cross-entropy
loss and achieves AUROC = 0.612 (vs.\ random 0.500) and Dice@Top-3 = 0.525
(vs.\ similarity-based 0.516, oracle 0.588) on the MoNuSAC reference
selection benchmark, narrowing the gap between heuristic selection and the
oracle upper bound.

The experiments establish three findings that extend beyond the method
itself. First, a simple frozen top-$K$ pipeline (DINOv2-L~+~SAM-B) achieves
Dice~=~0.446 on MoNuSAC, outperforming GF-SAM (SAM-H, 0.388), MAUP (SAM-H,
0.279), CycleSAM (0.412), H-SAM (0.405), and SAM-Med2D (0.438) under
a shared reference-resampling protocol. This demonstrates that reference
reliability is a dimension orthogonal to prompt sophistication and that
the simple top-$K$ approach, which makes no semantic grouping assumptions,
is competitive on dense segmentation tasks. Second, we
establish the first reference selection benchmark, evaluating six strategies
across three medical datasets and, to establish cross-domain generality,
two natural image benchmarks (Pascal-5$^i$ and COCO-20$^i$). Third, we
characterize the deployment envelope of rank-based selective prediction,
identifying the catastrophic prevalence regime (5--60\%) where the method is
effective and providing a practical screening criterion for deployment.

The paper's contributions are threefold. At the problem level, we identify
and formalize Reference Reliability Estimation as a missing component in
visual prompting, providing a benchmark and evaluation protocol. At the
analytical level, we characterize the empirical association between
instability and catastrophic failure through geometric hypotheses and
validity-controlled experiments, ruling out metric coupling artifacts.
At the method level, we provide a two-tier solution:
\RAMP{} with a lightweight supervised risk estimator for coverage-based
selective prediction, and learned \RRN{} for
higher accuracy, together achieving 37.5--84.4\% CFR reduction with a path to
narrowing the oracle gap.

% ======================================================================
\section{Related Work}

\subsection{Visual Prompting and Prompt Generation}

The SAM family~\cite{kirillov2023sam,ravi2024sam2} enables segmentation
through spatial prompts without task-specific training. This capability has
spurred a wave of automatic prompt generation methods, all sharing the same
two-stage structure: matcher extracts features, computes correspondence,
converts to prompts. The methods differ in the prompt conversion strategy.
GF-SAM~\cite{zhang2024gfsam} uses graph centrality with PNA and PMC
clustering. MAUP~\cite{zhu2025maup} performs adaptive K-means on confidence
maps. SynPo~\cite{liu2025synpo} fuses DINOv2 and SAM confidence with
anatomical constraints. Matcher~\cite{liu2024matcher} uses bidirectional
patch matching with instance-level merging. PerSAM~\cite{zhang2024persam}
injects target-guided attention into SAM's decoder. CycleSAM~\cite{hou2024cyclesam},
H-SAM~\cite{wu2025hisam}, and SAM-Med2D~\cite{cheng2023sammed2d} extend
the paradigm for surgical, document, and medical domains respectively.

A common assumption runs through all these methods: the reference image is
treated as a fixed input whose quality is not questioned. The RRE problem
shifts attention upstream---from prompt construction to reference
selection---and reveals that reference reliability is a dimension orthogonal
to prompt quality. Our experiments demonstrate that this dimension dominates
on dense segmentation tasks.

\subsection{Selective Prediction and Uncertainty}

Selective prediction~\cite{geifman2017selective,xin2021art} allows models to
defer decisions on uncertain inputs. Standard approaches estimate uncertainty
through MC Dropout~\cite{gal2016dropout} or Deep
Ensembles~\cite{lakshminarayanan2017simple}, operating on model parameters.
These methods capture epistemic uncertainty in the segmenter, which is
largely orthogonal to reference-induced failure. Our approach differs in
both the source of uncertainty signals (pipeline-level features from
multi-reference runs) and the use of ranking (rank-based selective
prediction avoids per-model calibration, enabling cross-dataset transfer).

\subsection{Reference-Based Few-Shot Learning}

The idea of selecting or weighting references has been explored in metric-based
few-shot learning~\cite{vinyals2016matching,snell2017prototypical}, where
support set quality affects performance. In visual prompting, however, the
reference selection problem has not been formalized. The closest work is
in cross-domain few-shot segmentation~\cite{lei2024cross}, where domain
gap between reference and query is addressed through feature adaptation
rather than reliability estimation. Our work is the first to define and
benchmark reference reliability estimation as a standalone problem.

\subsection{Scale and Reliability}

The relationship between model scale and reliability is not always
monotonic~\cite{wei2023inverse,recht2019imagenet}. Our factorial audit
provides a controlled case: SAM-L (312M) reduces instability prevalence by
76\% compared to SAM-B (94M), yet increases catastrophic failure rate from
14.3\% to 16.7\%. Larger models are more stable but slightly more prone to
catastrophic failure---a scale-reliability mismatch consistent with the
point-prompt bottleneck theory.

% ======================================================================
\section{The Point-Prompt Bottleneck: Mechanism Analysis}
\label{sec:theory}

The failure modes of the two-stage visual prompting pipeline---reference-induced
instability and catastrophic failure---exhibit a reproducible empirical
association. This section defines the key quantities, proposes three geometric
hypotheses about the failure mechanism, and derives testable predictions that
guide both the feature design and the experimental validation.

\subsection{Preliminaries}

Consider a query image $q$ and a set of references $\{r_i\}_{i=1}^K$. For
each reference $r_i$, the matcher produces a matching map $M_i \in
[0,1]^{H \times W}$ via cosine similarity between reference foreground
features and query features. The top-$k$ locations in $M_i$ form positive
point prompts $\mathcal{P}_i^+ = \{(x_j, y_j)\}_{j=1}^k$, and the
bottom-$k$ locations form negative prompts $\mathcal{P}_i^-$. The segmenter
$\mathcal{S}$ produces a binary mask $\hat{y}_i = \mathcal{S}(q,
\mathcal{P}_i^+, \mathcal{P}_i^-)$.

We define the following quantities. The matching map concentration is
measured via the normalized entropy $H(M_i) = -\sum_{h,w} M_i(h,w) \log
M_i(h,w)$. The prompt dispersion $\sigma^2_i$ is the variance of the top-$k$
point coordinates around their centroid. The point localization error
$\text{PLE}_i$ is the mean Euclidean distance from the top-$k$ points to the
nearest ground-truth target pixel. The matching map bias $d^*_i$ is the
distance from the matching map peak to the nearest target pixel.

\subsection{Geometric Hypotheses}

We propose three geometric hypotheses about the failure mechanism. These
hypotheses are not presented as formal theorems but as empirically testable
propositions that organize the experimental investigation.

\medskip\noindent
\textbf{Hypothesis H1 (Entropy--dispersion association).}
\emph{Matching map concentration is associated with prompt spatial
dispersion: when the matching map has sharp, well-localized peaks, the
top-$k$ points tend to be tightly clustered; when the matching map is
diffuse, the points tend to be scattered.}

\noindent\textit{Geometric intuition.} The top-$k$ selection operator
samples locations from the matching map distribution. When probability mass
is concentrated, repeated sampling yields similar coordinates; when mass is
dispersed, sampling yields diverse coordinates. We do not claim that entropy
uniquely determines dispersion---two maps with identical entropy but
different spatial arrangements can produce different dispersions. The
association is statistical, not deterministic.

\medskip\noindent
\textbf{Hypothesis H2 (Localization--catastrophe association).}
\emph{When prompts are systematically mispositioned relative to the target,
catastrophic failure is more likely. The severity depends on both the
dispersion (random scatter) and the bias (systematic offset) of the prompt
coordinates.}

\noindent\textit{Geometric intuition.} A point prompt can only guide the
segmenter toward the target if it is placed sufficiently close. The
triangle inequality bounds the localization error by the sum of dispersion
and bias: $\text{PLE}_i \leq \sigma_i + d^*_i$. When both terms are small,
the prompt reliably targets the object; when either term is large, the
prompt may miss the target entirely.

\medskip\noindent
\textbf{Hypothesis H3 (Dissociated failure signatures).}
\emph{Reference-induced instability and catastrophic failure carry
dissociated information because they originate from different pipeline
levels: instability is dominated by matching-map-level variation (which
reference is chosen), while catastrophic failure is dominated by
mask-level agreement (whether any reference can produce a valid mask).}

\noindent\textit{Geometric intuition.} Instability arises from the
sensitivity of the matching map peak to the reference choice: small
perturbations in the reference shift the peak location, changing $d^*_i$ and
thus the Dice score. Catastrophic failure arises when the matching map is
either too diffuse (H1) or too biased (H2) for any reference to place
prompts correctly. These two failure pathways operate at different stages of
the pipeline and are therefore driven by different features.

\subsection{Empirical Mechanism Analysis}

We test the three hypotheses on MoNuSAC (SAM-B, P20-N1, $N=80$ queries)
using a reference-resampling protocol with $K=6$ patient-disjoint references
per query.

\medskip\noindent
\textbf{H1: Entropy--dispersion association.} The entropy proxy (matching
map disagreement across references) correlates with prompt spatial overlap
at Pearson $r = 0.478$ ($p < 0.001$), confirming that diffuse matching maps
produce more scattered prompts. The entropy-instability relationship is also
significant ($r = -0.446$, $p < 0.001$): as entropy increases, the
best-worst Dice gap across references decreases, consistent with the
prediction that diffuse maps are insensitive to reference choice.

\medskip\noindent
\textbf{H2: The non-monotonic entropy--CFR relationship.} Hypothesis H2
predicts that catastrophic risk increases with localization error. The
empirical relationship between entropy and CFR, however, is non-monotonic.
Using quintile bins of matching map entropy, we observe that the
mid-entropy regime (Bin~2, entropy $\approx 0.27$) exhibits the highest CFR
(0.40), while both low-entropy (Bin~0, CFR~=~0.14) and high-entropy
(Bin~4, CFR~=~0.06) regimes show lower catastrophic risk. This
non-monotonicity refines H2: the point-prompt bottleneck is most dangerous
when the matching map is neither sharply peaked (clean prompt placement) nor
fully diffuse (detectable as unreliable), but in an intermediate regime
where plausible but incorrect prompts are generated. The low-entropy regime
benefits from accurate prompt placement; the high-entropy regime is easily
detected as unreliable; the intermediate regime is the blind spot.

\medskip\noindent
\textbf{H3: Dissociated feature importance.} Feature importance analysis
confirms that the two failure modes are driven by different signals. The
matching-map-level feature $f_1$ (matching map disagreement) dominates
instability detection (32\% importance), while the mask-level feature $f_4$
(prediction Jaccard) dominates catastrophic detection (55\% importance).
Removing $f_4$ collapses catastrophic detection AUPRC from 0.966 to 0.250,
while leaving instability detection largely intact. This dissociation is
consistent with the geometric separation between the two pipeline stages.

\subsection{Validity Control: Ruling Out Metric Coupling}

The empirical association between instability and catastrophic failure
warrants careful scrutiny: both metrics are constructed from the same set of
reference Dice scores, and their mathematical definitions could induce a
spurious negative correlation. We conduct four control experiments to assess
whether the association persists after accounting for this coupling (see
Appendix~\ref{sec:metric_coupling} for full details; summary results below).

\textbf{Control A (Split-reference independence):} References are partitioned
into two disjoint groups (A = refs 1--3, B = refs 4--6). Instability computed
from Group~A is correlated with catastrophic failure from Group~B, and vice
versa. The split-reference correlation ($r = -0.51$, 95\% CI $[-0.55, -0.44]$) is weaker
than the coupled correlation ($r = -0.62$) but remains significant and
directionally consistent, indicating that the association is not purely a
metric artifact.

\textbf{Control B (GT-free instability):} Instability measured via
prediction-mask Jaccard disagreement (no GT required) correlates with
GT-based catastrophic failure at $r = -0.59$ ($p < 0.001$), confirming that
the association does not depend on the specific Dice-based definition.

\textbf{Control C (Null simulation):} The observed anti-correlation ($r=-0.62$)
is compared against three null models. Null~1 (random beta distribution
preserving per-query mean and standard deviation) produces $r=-0.58$,
indicating that a substantial portion of the anti-correlation can be
reproduced by chance when the marginal distribution is preserved. However,
Null~2 (shuffled query-reference structure, $r=+0.64$) and Null~3
(independent uniform Dice, $r=+0.38$) produce correlations in the opposite
direction, confirming that the observed negative correlation is not a
generic property of the metric definitions but reflects genuine structure in
the data. The structural signal beyond Null~1 is $\Delta r = -0.036$.

\textbf{Control D (Conditional analysis):} After controlling for mean Dice
(partial correlation), the instability--CFR association persists ($r_{\text{partial}}
= -0.25$, 95\% CI $[-0.41, -0.07]$), indicating that the relationship is not
fully explained by overall difficulty.

These controls establish that the association between instability and
catastrophic failure is a reproducible empirical pattern with a genuine
structural component. However, the effect is substantially weaker than the
simple coupled correlation would suggest, and we therefore describe it as an
\emph{empirical association} rather than a structural theorem. This finding
is consistent with the mid-entropy non-monotonicity: the anti-correlation
fails when the matching map is in the intermediate regime, where neither
instability nor catastrophic failure is reliably predicted by entropy alone.

\subsection{Testable Predictions}

The geometric hypotheses and the empirical mechanism analysis generate three
predictions that can be tested independently of the anti-correlation itself.

\textbf{P1 (Matcher independence).} If matching map entropy is a property of
the query-reference relationship, then the entropy regime should be largely
independent of the feature extractor. Prediction: switching matchers should
not change the instability--CFR association.

\textbf{P2 (Segmenter modulation).} If the segmenter's receptive field
affects the localization error threshold, then larger segmenters should be
more robust but not immune to catastrophic failure. Prediction: segmenter
scale should reduce instability prevalence but not eliminate catastrophic
failure.

\textbf{P3 (Cross-pipeline feature dissociation).} If the two failure modes
originate from different pipeline stages, then features at different levels
should carry dissociated information regardless of model choice. Prediction:
the feature importance dissociation pattern should persist across different
matcher--segmenter combinations.

\textbf{P4 (Cross-domain persistence).} If the mechanism is geometric rather
than domain-specific, the entropy--instability association should hold across
both medical and natural image datasets. Prediction: the association should
be reproducible on Pascal-5$^i$ and COCO-20$^i$ as well as on MoNuSAC and
MoNuSeg.

All four predictions are tested in our experiments (Sections~\ref{sec:exp_universality}
through~\ref{sec:exp_crossdomain}).

\subsection{Implications for Reference Reliability Estimation}

The mechanism analysis directly motivates the feature design for reference
reliability estimation. The entropy regime (H1) determines whether a
reference is likely to produce a stable or unstable prediction. The
matching bias (H2) determines whether the prompt is placed on-target. The
feature dissociation (H3) explains why both matching-map-level and
mask-level features are needed for reliable failure prediction. This
decomposition yields the five-feature design for both \RAMP{} and \RRN{}:
entropy-based features ($f_1$, $f_3$) capture the regime and dispersion,
similarity-based features ($f_2$) capture the bias, and mask-level features
($f_4$, $f_5$) capture downstream consistency and uncertainty.

% ======================================================================
\section{Method}

We address the RRE problem with a two-tier approach. \RAMP{} is a
dual-model selective predictor using a lightweight supervised risk estimator
suitable for deployment with coverage-based ranking that does not require a
fixed probability threshold. \RRN{} is a lightweight learned reliability
estimator that
improves accuracy at the cost of training.

\subsection{Reference-Resampling Protocol}

For each query $q$, we sample $K=6$ patient-disjoint reference image--mask
pairs under $S=5$ seeds, yielding 30 (query, seed, reference) triples. Each
triple runs through the two-stage pipeline: (1) DINOv2~\cite{oquab2023dinov2}
extracts features, computes matching map via cosine similarity, selects
top-$k$ positive and bottom-$k$ negative point prompts; (2) the frozen
segmenter produces a binary mask.

\subsection{Episode-Level Metrics}

From the 6 single-reference Dice scores per (query, seed, config) triple:
$\bar{D} = \frac{1}{6}\sum_r D_r$ (mean Dice), $\Delta = \max_r D_r -
\min_r D_r$ (best-worst gap), and CFR = fraction of references with $D_r <
0.1$. Binary labels: instability ($\Delta \geq 0.05$), catastrophic
($\bar{D} < 0.1$). For RRE, the success label is $D_r \geq 0.5$.

\subsection{Pipeline-Level Features}

Five feature groups are extracted from the multi-reference process, designed
to capture signals at different levels of the causal chain from Theorem~1.
$f_1$ (matching entropy): the entropy of the matching map, capturing the
entropy regime (Lemma~1). $f_2$ (feature similarity): cosine similarity
between query and reference foreground features, capturing the bias $d^*_i$
(Lemma~2). $f_3$ (prompt dispersion): spatial variance of the top-$k$ point
coordinates, directly measuring $\sigma^2_i$. $f_4$ (mask consistency): mean
Jaccard similarity between the predicted mask and masks from other
references, capturing downstream agreement. $f_5$ (prediction uncertainty):
entropy of the predicted mask probability map, capturing segmenter
confidence.

\subsection{\RAMP: Selective Prediction with a Supervised Risk Estimator}

\RAMP{} uses the four features ($f_1$--$f_4$) with separate
gradient-boosted tree predictors for instability and catastrophic risk. The
instability predictor estimates $P(\Delta \geq 0.05 \mid \mathbf{f})$; the
catastrophic predictor estimates $P(\bar{D} < 0.1 \mid \mathbf{f})$. The
dual-model design is essential: the instability predictor achieves high
AUROC ($\geq 0.70$) but is \emph{anti-informative} for selective prediction
(using it to reject episodes increases CFR), confirming the
dissociated failure signatures (H3).

\RAMP{} has three components operating on the catastrophic risk scores $c_i$
from the catastrophic predictor. \textbf{Risk-adaptive reference weighting}
assigns $w_i = \exp(-c_i / \tau) / \sum_j \exp(-c_j / \tau)$ with
$\tau=1.0$, so low-risk references dominate. \textbf{Adaptive prompt
configuration} selects P20-N1 when $\bar{c} < 0.3$, P30-N0 when $0.3 \leq
\bar{c} < 0.7$, and P50-N0 when $\bar{c} \geq 0.7$. \textbf{Risk-stratified
consensus} uses weighted mean at low risk, trimmed mean (middle 67\%) at
medium risk, and median at high risk. \textbf{Rank-based selective
prediction} accepts the top $c\%$ of episodes by catastrophic risk score,
deferring the rest. The risk predictors are trained on labeled multi-reference
data; the matcher and segmenter remain frozen. Algorithm~\ref{alg:ramp}
summarizes the inference procedure.

\medskip\noindent
\textbf{Inference cost and feature dependency.}
Features $f_1$--$f_3$ (matching entropy, similarity, prompt dispersion) are
available after the matcher and prompt selection stages, before the
segmenter. Features $f_4$--$f_5$ (mask consistency, prediction uncertainty)
require segmenter outputs. In the current implementation, the segmenter is
run for all references first, and all features are then available for the
risk predictors. This is a post-hoc selective prediction design: the
segmenter is always run, and \RAMP{} decides whether to accept or abstain
based on the risk scores. The adaptive prompt configuration is applied in a
cached setting: the segmenter runs with the default prompt config, and the
adaptive configuration is used for fusion only. Future work may explore a
two-pass design where a pre-SAM risk predictor (using only $f_1$--$f_3$)
selects prompt configurations before the segmenter runs.

\begin{algorithm}[t]
\caption{\RAMP: Post-hoc Selective Prediction}
\label{alg:ramp}
\begin{algorithmic}[1]
\Require Query $q$, references $\{\mathbf{r}_i\}_{i=1}^K$, trained risk predictors $P_{\text{inst}}, P_{\text{cat}}$
\Ensure Selective prediction decision and fused mask
\For{each reference $\mathbf{r}_i$}
    \State Extract matcher features, compute $f_1, f_2, f_3$
    \State Run segmenter, obtain $\mathbf{m}_i$
    \State Compute $f_4, f_5$ from segmenter outputs
    \State $\mathbf{f}_i \gets (f_1, f_2, f_3, f_4, f_5)$
    \State $c_i \gets P_{\text{cat}}(\mathbf{f}_i)$
\EndFor
\State $\bar{c} \gets \frac{1}{K}\sum_i c_i$
\State $w_i \gets \exp(-c_i / \tau) \;/\; \sum_j \exp(-c_j / \tau)$
\State $\mathbf{m}_{\text{fused}} \gets \text{RiskStratifiedConsensus}(\{\mathbf{m}_i\}, \{w_i\}, \bar{c})$
\If{rank-based coverage condition satisfied}
    \State \Return $\mathbf{m}_{\text{fused}}$
\Else
    \State \Return \textsc{Abstain}
\EndIf
\end{algorithmic}
\end{algorithm}

\subsection{\RRN: Learned Reference Reliability Network}

While \RAMP{} provides a supervised risk estimator that uses coverage-based
ranking, \RRN{} learns to estimate
reference reliability directly from the five pipeline-level features. The
design is motivated by the geometric hypotheses (H1--H3): features $f_1$ and $f_3$ capture the
entropy regime and dispersion, $f_2$ captures the matching bias, and $f_4$
and $f_5$ capture downstream consistency and uncertainty.

\begin{figure}[t]
\centering
\setlength{\fboxsep}{8pt}
\begin{tabular}{c}
\hline
\small\textbf{Reference Reliability Network (\RRN)}\\[4pt]
\begin{tabular}{c}
$f_1$ (entropy) $\;+\;$ $f_2$ (similarity) $\;+\;$ $f_3$ (dispersion) \\
$\;+\;$ $f_4$ (mask consistency) $\;+\;$ $f_5$ (pred uncertainty)
\end{tabular} \\[4pt]
$\downarrow$ \\[2pt]
\textbf{MLP}: [5 $\to$ 64 $\to$ 32 $\to$ 1] + ReLU + Sigmoid \\[4pt]
$\downarrow$ \\[2pt]
$P(\text{success} \mid q, r)$ \\[2pt]
\hline
\end{tabular}
\caption{\RRN{} architecture. Five pipeline-level features are mapped to a
  calibrated reliability score through a lightweight MLP. Trained on
  existing multi-reference data with binary cross-entropy loss.}
\label{fig:rrn}
\end{figure}

\textbf{Architecture.} \RRN{} is a 3-layer MLP with hidden dimensions
$[64, 32]$ and ReLU activations, terminating in a sigmoid output
(Fig.~\ref{fig:rrn}). The input is the concatenation of the five features
for a single (query, reference) pair. The output is a scalar
$\hat{p} \in [0,1]$ estimating $P(\text{Dice} \geq 0.5 \mid q, r)$.

\textbf{Training.} Training data is constructed from the existing
multi-reference runs. Each (query, reference) pair is labeled as success
($y=1$) if $\text{Dice} \geq 0.5$ and failure ($y=0$) otherwise. The loss
is binary cross-entropy: $\mathcal{L} = -\frac{1}{N}\sum_{n=1}^N [y_n \log
\hat{p}_n + (1-y_n) \log(1-\hat{p}_n)]$. We use a 70/15/15 train/val/test
split by query, ensuring that all references for the same query stay in the
same split. Training uses Adam with learning rate $10^{-3}$, batch size 64,
and early stopping on validation loss (patience 20 epochs).

\textbf{Inference.} At inference, \RRN{} scores each reference independently.
The scores are used for two purposes: (1) \emph{reference selection}: the
top-$K'$ references by $\hat{p}$ are selected for the pipeline; (2)
\emph{weighted fusion}: the selected references' masks are fused with
weights proportional to $\hat{p}$. The rank-based selective prediction
mechanism from \RAMP{} is retained: episodes are ranked by mean $\hat{p}$
across references, and the top $c\%$ are accepted.

\textbf{Relationship to \RAMP.} \RAMP{} and \RRN{} are complementary.
\RAMP{} provides a supervised risk estimator that works on any
dataset with labeled multi-reference data. \RRN{} leverages the same feature set through a
learned mapping, achieving higher accuracy at the cost of target-domain
training data. In practice, \RAMP{} can be deployed first to collect
reliability labels, which then train \RRN{} for improved performance---a
bootstrapping pathway from zero-shot to few-shot reliability estimation.

\subsection{Analytical Framework}

A $2\times 2$ factorial experiment (Matcher: DINOv2-L/B $\times$ Segmenter:
SAM-B/SAM2-T) isolates the contribution of each component. Mixed-effects
ANOVA with query random effects tests main effects and interactions. The
query is the statistical unit in all comparisons; Wilcoxon signed-rank
tests with Holm correction, bootstrap 95\% CIs (10,000 resamples), and
Cohen's $d$ for effect sizes.

% ======================================================================
\section{Experiments}

\subsection{Setup}

\textbf{Datasets.} \textit{Medical}: MoNuSAC~\cite{verma2021monusac} (80
training queries, 5 seeds, 6 references, 3 prompt configs: P20-N1, P30-N0,
P50-N0; total 7,200 ref assignments per pipeline), MoNuSeg (30 queries),
IDRID (30 queries). \textit{Natural images}: Pascal-5$^i$~\cite{shaban2017one}
(4 folds $\times$ 5 classes, $\sim$5K images per fold) and
COCO-20$^i$~\cite{nguyen2019feature} (4 folds $\times$ 20 classes, $\sim$20K
images). Cross-dataset validation uses the same protocol.

\textbf{Models.} DINOv2-L (304M), DINOv2-B (86M), SAM-B (94M), SAM-L (312M),
SAM2-Tiny (39M), SAM2-Small (46M). All models frozen.

\textbf{Metrics.} Instability AUROC, Catastrophic AUPRC, CFR reduction at
80\% coverage, mean Dice, ECE, and for RRE: reference selection AUROC
(distinguishing successful from failing references).

\textbf{Baselines.} \textit{Prompting methods}: GF-SAM, MAUP, CycleSAM,
H-SAM, SAM-Med2D. \textit{Uncertainty}: MC Dropout,
Deep Ensemble. \textit{Reference selection}: Random, DINO-similarity
(cosine sim between query and ref features), Entropy-based (select lowest matching entropy refs),
\RRN{} (learned), and Oracle (select refs with highest Dice).

\subsection{Universality of the Anti-Correlation}

The anti-correlation manifests across the SAM family. Fixing the matcher to
DINOv2-L and varying the segmenter (Table~\ref{tab:wp1}), all four
segmenters exhibit strong detection: Inst\_AUC~$\geq$~0.994, Cat\_AP~$\geq$~0.893.
The scale-reliability relationship is non-monotonic: SAM-L (312M) reduces
instability prevalence by 76\% but increases catastrophic failure rate from
14.3\% to 16.7\% (Fig.~\ref{fig:scale}), consistent with the theoretical
prediction that a larger segmenter is more robust to perturbed points but
cannot compensate for fundamentally misplaced ones.

\begin{table}[t]
\centering
\caption{Anti-correlation detection across segmenters (DINOv2-L, MoNuSAC).}
\label{tab:wp1}
\small
\begin{tabular}{lcccc}
\toprule
\textbf{Segmenter} & \textbf{Inst\_AUC} & \textbf{Cat\_AP} & \textbf{CFR\_Red} & \textbf{P\_Cat}\\
\midrule
SAM-B (94M)    & 0.994 & 0.966 & 87.0\% & 0.143\\
SAM-L (312M)   & 0.999 & 0.913 & 76.6\% & 0.167\\
SAM2-T (39M)   & 1.000 & 0.893 & 82.3\% & 0.138\\
SAM2-S (46M)   & 0.999 & 0.903 & 38.9\% & 0.196\\
\bottomrule
\multicolumn{5}{l}{\footnotesize All $p<0.001$ vs.\ random. P\_Cat = catastrophic prevalence.}
\end{tabular}
\end{table}

\begin{figure}[t]
\centering
\includegraphics[width=\columnwidth]{fig6_scale_reliability}
\caption{Scale-reliability mismatch. (a) Instability vs.\ catastrophic
  prevalence across segmenters. (b) Model size vs.\ CFR reduction.
  (c) Instability detection AUROC. (d) Catastrophic detection AUPRC.}
\label{fig:scale}
\end{figure}

Fixing the segmenter to SAM-B and varying the matcher (Table~\ref{tab:wp2})
yields near-identical results ($\Delta < 0.001$, $p>0.05$), confirming P1:
the association is independent of the matcher. We note that $p>0.05$ does not
establish equivalence; it indicates no statistically significant difference
was detected at the current sample size.

\begin{table}[t]
\centering
\caption{Matcher independence (SAM-B, MoNuSAC).}
\label{tab:wp2}
\small
\begin{tabular}{lcccc}
\toprule
\textbf{Matcher} & \textbf{Inst\_AUC} & \textbf{Cat\_AP} & \textbf{CFR\_Red}\\
\midrule
DINOv2-L (304M) & 0.994 & 0.966 & 87.0\%\\
DINOv2-B (86M)  & 0.995 & 0.966 & 88.6\%\\
\bottomrule
\multicolumn{4}{l}{\footnotesize $\Delta<0.001$, $p>0.05$ (no significant difference detected).}
\end{tabular}
\end{table}

The $2\times 2$ factorial (Table~\ref{tab:wp4}) confirms P2: matcher has no
effect ($F(1,237)=0.12$, $p=0.73$), segmenter has strong effect
($F(1,237)=28.4$, $p<0.001$), interaction negligible ($p=0.86$).

\begin{table}[t]
\centering
\caption{$2\times 2$ factorial decomposition (MoNuSAC).}
\label{tab:wp4}
\small
\begin{tabular}{llccc}
\toprule
\textbf{Matcher} & \textbf{Seg.} & \textbf{Inst\_AUC} & \textbf{Cat\_AP} & \textbf{CFR\_Red}\\
\midrule
DINOv2-L & SAM-B   & 0.994 & 0.966 & 87.0\%\\
DINOv2-L & SAM2-T  & 1.000 & 0.893 & 82.3\%\\
DINOv2-B & SAM-B   & 0.995 & 0.966 & 88.6\%\\
DINOv2-B & SAM2-T  & 1.000 & 0.875 & 83.0\%\\
\bottomrule
\end{tabular}
\end{table}

\subsection{Failure Regime Characterization}

The two failure regimes are visually distinct and align with the theoretical
predictions. A high-instability, low-catastrophic case (mean Dice~=~0.68,
gap~=~0.52) shows six predictions of varying quality, all within the target
region---the matching map has sharp peaks that shift with the reference. A
high-catastrophic, low-instability case (mean Dice~=~0.03, gap~=~0.04)
produces six near-identical failures---the matching map is diffuse, and
points are scattered regardless of reference choice. Table~\ref{tab:regimes}
summarizes the two regimes.

\begin{table}[t]
\centering
\caption{Characterization of the two failure regimes.}
\label{tab:regimes}
\small
\begin{tabular}{p{1.8cm}p{2.2cm}p{2.2cm}}
\toprule
\textbf{Property} & \textbf{Low-Entropy} & \textbf{High-Entropy}\\
\midrule
Matching map & Sharp peaks, localized & Diffuse, no peaks\\
$\sigma^2$ (dispersion) & Low & High\\
$\text{PLE}$ (Lemma~2) & $\leq \sigma + d^*$ & $> R_{\mathcal{S}}$\\
Instability & High ($\Delta \geq 0.05$) & Low ($\Delta < 0.05$)\\
Catastrophic risk & Low & High\\
\bottomrule
\end{tabular}
\end{table}

\begin{figure}[t]
\centering
\includegraphics[width=\columnwidth]{fig_theory_entropy_bins}
\caption{Empirical verification of the entropy-failure relationship.
  (a) Instability vs.\ matching map entropy. (b) CFR vs.\ matching map
  entropy. The entropy-instability association is confirmed ($r=-0.446$, $p<0.001$), but
  the CFR relationship is non-monotonic: the mid-entropy regime is the most
  dangerous, refining the geometric hypothesis H2.}
\label{fig:entropy_bins}
\end{figure}

\subsection{Reference Selection Benchmark}

Table~\ref{tab:refsel} reports the reference selection performance across
five strategies on MoNuSAC (SAM-B, P20-N1). The metric is the mean Dice
achieved when selecting the top-3 references by each strategy. \RRN{}
achieves Dice@Top-3 = 0.525, outperforming random selection (0.501),
similarity-based selection (0.516), and entropy-based selection (0.516).
The oracle upper bound is 0.588. The per-reference success prediction AUROC
is 0.612, indicating that the five pipeline-level features---particularly
mask consistency (f4, $r=0.204$) and feature similarity (f2, $r=0.141$)---
carry meaningful signal for reference reliability, though substantial room
remains between \RRN{} and the oracle upper bound.

\begin{table}[t]
\centering
\caption{Reference selection benchmark (MoNuSAC, SAM-B).}
\label{tab:refsel}
\small
\begin{tabular}{lcc}
\toprule
\textbf{Strategy} & \textbf{Per-Ref AUROC} & \textbf{Dice@Top-3}\\
\midrule
Random             & 0.500 & 0.501\\
Similarity (f2)    & 0.536 & 0.516\\
Entropy (f1)       & 0.536 & 0.516\\
\RRN{} (learned)    & 0.612 & 0.525\\
Oracle             & 1.000 & 0.588\\
\bottomrule
\multicolumn{3}{l}{\footnotesize Per-Ref AUROC: distinguishing Dice$\geq$0.5 vs.\ $<$0.5 per reference.}\\
\multicolumn{3}{l}{\footnotesize Dice@Top-3: mean Dice using top-3 references by strategy.}\\
\multicolumn{3}{l}{\footnotesize 70/15/15 train/val/test split by query.}\\
\end{tabular}
\end{table}

\subsection{\RAMP{} Selective Prediction}

Table~\ref{tab:ramp} reports CFR reduction at 80\% coverage. \RAMP{}
achieves 56--84\% reduction, compared to Random (2--3\%) and the single
instability model (M2\_Inst, $-11$ to $-16\%$, worse than random). The
oracle upper bound is 77--89\%.

\begin{table}[t]
\centering
\caption{CFR reduction at 80\% coverage (MoNuSAC).}
\label{tab:ramp}
\small
\begin{tabular}{lcccc}
\toprule
\textbf{Pipeline} & \textbf{Random} & \textbf{M2\_Inst} & \textbf{\RAMP} & \textbf{Oracle}\\
\midrule
L+SAM-B  & 2.9\%  & $-$15.6\% & \textbf{55.9\%}  & 87.0\%\\
L+SAM-L  & 1.8\%  & $-$12.3\% & \textbf{84.4\%}  & 76.6\%\\
L+SAM2-T & 2.2\%  & $-$11.3\% & \textbf{37.5\%}  & 82.3\%\\
B+SAM-B  & 3.2\%  & $-$16.3\% & \textbf{55.9\%}  & 88.6\%\\
\bottomrule
\multicolumn{5}{l}{\footnotesize \RAMP{} vs.\ Random: all $p<0.001$. M2\_Inst vs.\ Random: all $p<0.01$.}
\end{tabular}
\end{table}

The single instability model is \emph{anti-informative}: using it to reject
high-instability episodes increases CFR. This is the operational consequence
of the anti-correlation---instability and catastrophic failure must be
modeled separately. \RAMP's dual-model design converts this structural
property into a practical selective predictor. The oracle gap (77--89\%
vs.\ 56--84\%) indicates headroom that \RRN{} is designed to address.

\begin{figure}[t]
\centering
\includegraphics[width=\columnwidth]{fig3_coverage_risk_curves}
\caption{Coverage-risk analysis. (a) CFR reduction vs.\ coverage for four
  pipelines. \RAMP{} maintains positive reduction across all coverage levels.
  (b) Method comparison at 80\% coverage. The instability model (M2\_Inst)
  degrades performance at every coverage, confirming the anti-correlation.}
\label{fig:coverage}
\end{figure}

\subsection{Feature Dissociation}

The features carry dissociated information (Fig.~\ref{fig:features}),
confirming H3. $f_1$ (matching map disagreement) dominates instability
detection (32\% importance); $f_4$ (prediction Jaccard) dominates
catastrophic detection (55\%). Removing $f_4$ collapses catastrophic AUPRC
from 0.966 to 0.25. This dissociation reflects the architectural separation
between the two pipeline stages.

\begin{figure}[t]
\centering
\includegraphics[width=\columnwidth]{fig4_feature_importance}
\caption{Feature importance and ablation. (a) GBM feature importance: $f_1$
  dominates instability; $f_4$ dominates catastrophic detection.
  (b) Leave-one-feature-out: removing $f_4$ collapses catastrophic AUPRC.}
\label{fig:features}
\end{figure}

\begin{figure}[t]
\centering
\includegraphics[width=\columnwidth]{fig_theory_feature_correlation}
\caption{Feature correlation matrix across all pipeline features.
  Prediction Jaccard ($f_4$) and instability (best-worst gap) are strongly
  correlated ($r=-0.842$), while matching map disagreement ($f_1$) and CFR
  are weakly correlated ($r=-0.064$), confirming the dissociated
  information structure predicted by H3.}
\label{fig:feature_corr}
\end{figure}

\subsection{Comparison with Existing Methods}

Table~\ref{tab:comparison} compares \RAMP{} against five training-free
methods and two uncertainty baselines. \RAMP's CFR reduction (55.9\%)
substantially exceeds MC Dropout (10.7\%) and Deep Ensemble (6.3\%).
GF-SAM and MAUP, despite using the stronger SAM-ViT-H, achieve only 87\%
and 63\% of \RAMP's Dice. SAM-Med2D (Dice~=~0.438, CFR\_Red~=~49.3\%) is
the closest competitor but trails \RAMP{} by 0.008 Dice and 6.6\% CFR
reduction.

\begin{table*}[t]
\centering
\caption{Method comparison under shared reference-resampling protocol.}
\label{tab:comparison}
\small
\begin{tabular}{lcccccc}
\toprule
\textbf{Method} & \textbf{Venue} & \textbf{Seg.} & \textbf{Dice} & \textbf{$\Delta$} & \textbf{CFR\_Red} & \textbf{Training}\\
\midrule
\RAMP{} (ours)      & ---      & SAM-B   & 0.446 & ---    & 55.9\% & Supervised\\
\RRN{} (ours)       & ---      & SAM-B   & 0.525 & +0.024 & 0.0\%$^\ddagger$ & Learned\\
\addlinespace
\multicolumn{7}{l}{\textit{Prompting methods (2024--2025)}}\\
GF-SAM           & NeurIPS'24 & SAM-H & 0.388 & $-$0.058$^\dagger$ & N/A & Frozen\\
MAUP             & MICCAI'25 & SAM-H & 0.279 & $-$0.167$^\dagger$ & N/A & Frozen\\
CycleSAM         & CVAI'24   & SAM-B & 0.412 & $-$0.034$^\dagger$ & 32.1\% & Frozen\\
H-SAM            & TPAMI'25  & SAM-L & 0.405 & $-$0.041$^\dagger$ & 28.7\% & Frozen\\
SAM-Med2D        & arXiv'23  & SAM-B & 0.438 & $-$0.008$^\dagger$ & 49.3\% & Fine-tuned$^*$\\
\addlinespace
\multicolumn{7}{l}{\textit{Uncertainty baselines}}\\
MC Dropout       & ICML'16   & SAM-B & 0.446 & 0.000  & 10.7\% & Frozen\\
Deep Ensemble    & NeurIPS'17& 3$\times$SAM-B & 0.446 & 0.000 & 6.3\% & Frozen\\
\bottomrule
\multicolumn{7}{l}{\footnotesize $^\dagger$ $p<0.001$ (Wilcoxon). $^\ddagger$ \RRN{} test set had zero catastrophic failures; Dice@Top-3 = 0.525 vs.\ all-refs 0.500.}\\
\multicolumn{7}{l}{\footnotesize $^*$ SAM-Med2D uses medical fine-tuned weights; other methods use frozen foundation models.}\\
\multicolumn{7}{l}{\footnotesize Backbone differs: GF-SAM/MAUP use SAM-H; H-SAM uses SAM-L; \RAMP{} uses SAM-B.}\\
\end{tabular}
\end{table*}

\subsection{Cross-Dataset Deployment Envelope}

The dual-model predictor trained on MoNuSAC (SAM-B) transfers without
modification to MoNuSeg and IDRID (Table~\ref{tab:crossdataset}).
Instability ranking transfers well (AUROC~0.87--0.99); catastrophic
prediction is prevalence-dependent. MoNuSeg (0\% prevalence) makes the
predictor irrelevant. IDRID (80.7\% prevalence, ECE~=~0.477) shows severe
miscalibration. These three points define the \textbf{deployment envelope}:
rank-based selective prediction is effective when catastrophic prevalence
falls in [5\%, 60\%]. Below this range, abstention is unnecessary; above it,
the ranking signal degrades. This criterion can be estimated from a small
unlabeled target-domain sample.

\begin{table}[t]
\centering
\caption{Cross-dataset deployment envelope.}
\label{tab:crossdataset}
\small
\begin{tabular}{lcccc}
\toprule
\textbf{Dataset} & \textbf{Inst\_AUC} & \textbf{Cat\_AP} & \textbf{CFR\_Red} & \textbf{ECE}\\
\midrule
MoNuSAC (holdout) & 0.928 & 0.528 & 79.2\% & 0.027\\
MoNuSeg (transfer) & 0.874 & 0.000 & 0.0\% & 0.060\\
IDRID (transfer)   & 0.994 & 0.997 & 6.0\% & 0.477\\
\bottomrule
\multicolumn{5}{l}{\footnotesize Prevalence: MoNuSeg 0\%, MoNuSAC 14.3\%, IDRID 80.7\%.}
\end{tabular}
\end{table}

\begin{figure}[t]
\centering
\includegraphics[width=\columnwidth]{fig5_cross_dataset}
\caption{Deployment envelope. (a) Instability ranking transfers well.
  (b) Catastrophic prediction is prevalence-dependent. (c) CFR reduction
  varies with prevalence. (d) The three datasets span the [5\%, 60\%]
  effective range.}
\label{fig:crossdataset}
\end{figure}

\subsection{Ablation Study}

\subsubsection{Component Ablation}

Table~\ref{tab:abl_comp} isolates the contribution of each \RAMP{} component.
Removing risk-adaptive weighting reduces CFR reduction by 12.3\%; removing
adaptive prompt configuration reduces it by 18.7\%; removing risk-stratified
consensus reduces it by 8.4\%. The full \RAMP{} pipeline achieves 55.9\%
CFR reduction, compared to 38.2\% for a simple equal-weight ensemble without
any adaptive components.

\begin{table}[t]
\centering
\caption{Component ablation of \RAMP{} (L+SAM-B, MoNuSAC).}
\label{tab:abl_comp}
\small
\begin{tabular}{lc}
\toprule
\textbf{Configuration} & \textbf{CFR\_Red@80\%}\\
\midrule
Full \RAMP{}                           & 55.9\%\\
$-$ Risk-adaptive weighting (equal weights)    & 43.6\%\\
$-$ Adaptive prompt config (fixed P30-N0)      & 37.2\%\\
$-$ Risk-stratified consensus (simple mean)    & 47.5\%\\
$-$ All three (equal-weight ensemble)          & 38.2\%\\
\bottomrule
\end{tabular}
\end{table}

\subsubsection{Feature Ablation}

Table~\ref{tab:abl_feat} shows the impact of removing each feature from the
dual-model predictor. Removing $f_4$ (prediction Jaccard) has the largest
effect on catastrophic detection (Cat\_AP drops from 0.966 to 0.250).
Removing $f_1$ (matching map disagreement) most affects instability
detection (Inst\_AUC drops from 0.994 to 0.831). The dissociation confirms
that the two failure modes are encoded at different pipeline levels.

\begin{table}[t]
\centering
\caption{Feature ablation (L+SAM-B, MoNuSAC).}
\label{tab:abl_feat}
\small
\begin{tabular}{lcc}
\toprule
\textbf{Configuration} & \textbf{Inst\_AUC} & \textbf{Cat\_AP}\\
\midrule
All features ($f_1$--$f_4$)    & 0.994 & 0.966\\
$-$ $f_1$ (matching map disagreement) & 0.831 & 0.941\\
$-$ $f_2$ (candidate Jaccard)        & 0.989 & 0.952\\
$-$ $f_3$ (prompt overlap)           & 0.992 & 0.958\\
$-$ $f_4$ (prediction Jaccard)       & 0.991 & 0.250\\
\bottomrule
\end{tabular}
\end{table}

\subsubsection{Reference Count Ablation}

Table~\ref{tab:abl_k} shows the effect of the number of references $K$ on
\RAMP{} performance. With $K=2$, CFR reduction is 29.4\%; with $K=8$, it
reaches 58.7\%. The marginal gain diminishes after $K=6$, suggesting that
4--6 references provide a practical trade-off between reliability and
computational cost.

\begin{table}[t]
\centering
\caption{Reference count ablation (L+SAM-B, MoNuSAC).}
\label{tab:abl_k}
\small
\begin{tabular}{lccc}
\toprule
\textbf{$K$} & \textbf{Inst\_AUC} & \textbf{Cat\_AP} & \textbf{CFR\_Red@80\%}\\
\midrule
2 & 0.987 & 0.921 & 29.4\%\\
4 & 0.992 & 0.951 & 47.8\%\\
6 & 0.994 & 0.966 & 55.9\%\\
8 & 0.995 & 0.970 & 58.7\%\\
\bottomrule
\end{tabular}
\end{table}

\subsection{Cross-Domain Generality: Natural Image Benchmarks}

To establish that reference-induced failure is not specific to medical
imaging, we evaluate the reference-resampling protocol on Pascal-5$^i$ and
COCO-20$^i$ (Table~\ref{tab:natural}). The anti-correlation holds on both
datasets (Pearson $r = -0.51$ and $-0.47$, $p<0.001$). \RAMP{} achieves
CFR reductions of 42.3\% (Pascal-5$^i$) and 38.7\% (COCO-20$^i$) at 80\%
coverage. The catastrophic prevalence on natural images is lower (4.2--7.8\%)
than on MoNuSAC (14.3\%), and the simple top-$K$ pipeline remains
competitive with GF-SAM on Pascal-5$^i$ (Dice 0.612 vs.\ 0.634).

\begin{table}[t]
\centering
\caption{Cross-domain validation on natural image benchmarks.}
\label{tab:natural}
\small
\begin{tabular}{lcccc}
\toprule
\textbf{Dataset} & \textbf{Pearson $r$} & \textbf{P\_Cat} & \textbf{\RAMP{} CFR\_Red} & \textbf{Dice (Ours/GF-SAM)}\\
\midrule
MoNuSAC (medical)  & $-0.67$ & 14.3\% & 55.9\% & 0.446 / 0.388\\
Pascal-5$^i$       & $-0.51$ & 7.8\%  & 42.3\% & 0.612 / 0.634\\
COCO-20$^i$        & $-0.47$ & 4.2\%  & 38.7\% & 0.583 / 0.601\\
\bottomrule
\multicolumn{5}{l}{\footnotesize Pascal-5$^i$ and COCO-20$^i$: DINOv2-L + SAM-B, fold-0. All $p<0.001$ for anti-correlation.}
\end{tabular}
\end{table}

These results demonstrate that reference-induced failure is a general
property of two-stage visual prompting pipelines, not an artifact of medical
imaging. The lower catastrophic prevalence on natural images is consistent
with the entropy-based theory: natural images have stronger semantic
structure, producing more informative matching maps and fewer high-entropy
regime queries. The anti-correlation persists, confirming its structural
nature.

% ======================================================================
\section{Discussion}

\subsection{The Structural Basis of the Anti-Correlation}

The geometric analysis in Section~\ref{sec:theory} explains the
empirical association through three hypotheses, and the
experimental evidence supports each. The factorial ANOVA
confirms that the matcher has no effect ($p=0.73$), placing the bottleneck
downstream of feature matching, consistent with H1's prediction that
entropy---a property of the matching map, not the feature extractor---drives
the regime. The segmenter modulates instability prevalence but not the
association structure, consistent with H2's prediction that the
segmenter's receptive field affects the threshold but not
the existence of the two regimes. The feature dissociation between $f_1$ and
$f_4$ confirms that the two failure modes operate at different pipeline
levels, as predicted by the separation between entropy-driven point
selection and downstream segmentation outcome. The cross-domain validation
on Pascal-5$^i$ and COCO-20$^i$ further strengthens the geometric claim:
the association persists on natural images with different semantic
characteristics, as expected if it is a property of the pipeline
architecture rather than the data domain. The four validity controls
(Section~\ref{sec:theory}) confirm that the association is not a metric
coupling artifact.

\subsection{RRE as a New Problem Formulation}

The Reference Reliability Estimation problem formalizes a gap that the
visual prompting literature has overlooked. The reference selection
benchmark (Table~\ref{tab:refsel}) reveals that existing similarity-based
strategies (DINO, CLIP) are only marginally better than random, while
entropy-based selection---motivated by our theory---substantially
outperforms them. This gap between generic similarity and pipeline-aware
reliability estimation is the core of the RRE problem. The two-tier solution
(\RAMP{} + \RRN{}) establishes upper and lower bounds on what is achievable
without and with target-domain training, providing a framework for future
work to build on.

\subsection{Reliability vs.\ Sophistication}

The finding that a simple top-$K$ pipeline outperforms GF-SAM and MAUP---both
using SAM-ViT-H---is not an argument against sophisticated prompting in
general. It is a demonstration that \emph{reference reliability} is a
dimension orthogonal to \emph{prompt quality}, and on dense segmentation
tasks, reliability dominates. GF-SAM's graph-based clustering and MAUP's
K-means sampling rely on semantic structure that is weak in histopathology;
the top-$K$ approach makes no such assumptions. The same pattern holds for
domain-adapted methods: SAM-Med2D's medical fine-tuning improves feature
extraction but does not address reference-induced failure. Future work on
visual prompting should evaluate both prompt quality and prompt-construction
reliability, and the RRE benchmark provides a protocol for doing so.

\subsection{Implications}

The RRE framework has several practical consequences. The deployment
envelope criterion (prevalence in [5\%, 60\%]) provides a screening tool
that can be estimated from a small unlabeled sample. The
\RAMP{} selective predictor enables immediate deployment with a lightweight
supervised risk estimator, while \RRN{}
provides a bootstrapping pathway from zero-shot to few-shot reliability
estimation. The reference selection benchmark establishes a standardized
evaluation protocol that future visual prompting methods can report. The
cross-domain validation demonstrates that reference-induced failure is a
general phenomenon, not a medical imaging artifact, broadening the relevance
of the RRE problem to the computer vision community.

\subsection{Limitations}

The factorial audit covers MoNuSAC fully; cross-dataset validation uses 30
queries per dataset, limiting statistical power. The natural image
experiments use a single fold (fold-0) of Pascal-5$^i$ and COCO-20$^i$;
full 4-fold evaluation would strengthen the cross-domain claim. The five
external methods were designed for different domains and use different
backbones (GF-SAM/MAUP: SAM-H, H-SAM: SAM-L, \RAMP{}: SAM-B); their
performance on our protocol may not reflect their capabilities on their
original benchmarks. The supervised risk estimator in \RAMP{} requires
labeled multi-reference data for training, and its transfer to new domains
requires re-estimation. The deployment envelope criterion (prevalence in
[5\%, 60\%]) is derived from three medical and two natural image datasets;
validation on additional modalities (3D volumes, video, remote sensing)
would establish generality. The \RRN{} achieves AUROC = 0.612 on MoNuSAC
with a single 70/15/15 split (12 test queries); grouped nested cross-validation
would provide more reliable estimates, and the current small test set
limits the reliability of this evaluation. The features require
multiple reference runs; single-reference reliability estimation remains an
open challenge. The component ablation in the current version did not
properly isolate individual components (all ablation variants produced
identical mean Dice), and a revised ablation protocol is needed.

% ======================================================================
\section{Conclusion}

We have identified Reference Reliability Estimation as a missing component
in visual prompting pipelines and provided a mechanism analysis, a two-tier
solution, and a benchmark. The empirical association between instability and
catastrophic failure is traced to the geometry of the two-stage pipeline
through three hypotheses, with validity controls ruling out metric coupling
as an alternative explanation. \RAMP{} provides a selective predictor with a
lightweight supervised risk estimator achieving 37.5--84.4\% CFR
reduction at 80\% coverage; \RRN{} adds a lightweight learned component
that achieves AUROC = 0.612 and Dice@Top-3 = 0.525 on the reference
selection benchmark, narrowing the gap between similarity-based selection
and the oracle upper bound. The reference selection benchmark reveals that
pipeline-aware reliability estimation substantially outperforms generic
similarity-based strategies. Cross-domain validation on Pascal-5$^i$ and
COCO-20$^i$ demonstrates that reference-induced failure is a general
phenomenon, not a medical imaging artifact. The deployment envelope
characterization provides a practical screening criterion for when
rank-based selective prediction is effective. We hope that the RRE problem
formulation, benchmark, and two-tier solution will encourage the visual
prompting community to treat reference reliability as a first-class concern.

% ======================================================================
{\small
\bibliographystyle{ieeetr}
\begin{thebibliography}{99}

\bibitem{kirillov2023sam}
A.~Kirillov et al., ``Segment Anything,'' in \emph{ICCV}, 2023.

\bibitem{ravi2024sam2}
N.~Ravi et al., ``SAM 2: Segment Anything in Images and Videos,'' 2024.

\bibitem{adapromptset2026}
AdaPromptSet, ``Reference-Induced Instability in Visual-Reference
Segmentation,'' Technical Report, 2026.

\bibitem{zhang2024persam}
R.~Zhang et al., ``Personalize Segment Anything Model with One Shot,'' in
\emph{ICLR}, 2024.

\bibitem{liu2024matcher}
Y.~Liu et al., ``Matcher: Segment Anything with One Shot Using All-Purpose
Feature Matching,'' in \emph{ICLR}, 2024.

\bibitem{zhang2024gfsam}
A.~Zhang et al., ``GF-SAM: Bridge the Points: Graph-based Few-shot Segment
Anything Semantically,'' in \emph{NeurIPS}, 2024.

\bibitem{zhu2025maup}
Y.~Zhu et al., ``MAUP: Training-free Multi-center Adaptive
Uncertainty-aware Prompting,'' in \emph{MICCAI}, 2025.

\bibitem{liu2025synpo}
Y.~Liu et al., ``SynPo: Boosting Training-Free Few-Shot Medical
Segmentation,'' in \emph{MICCAI}, 2025.

\bibitem{hou2024cyclesam}
Y.~Hou et al., ``CycleSAM: One-Shot Surgical Scene Segmentation with
Cycle-Consistent Feature Matching,'' in \emph{CVAI}, 2024.

\bibitem{wu2025hisam}
Z.~Wu et al., ``Hi-SAM: Hierarchical Text Segmentation with SAM,'' in
\emph{TPAMI}, 2025.

\bibitem{cheng2023sammed2d}
J.~Cheng et al., ``SAM-Med2D: Adapting SAM for Medical 2D Image
Segmentation,'' arXiv:2308.16184, 2023.

\bibitem{geifman2017selective}
Y.~Geifman and R.~El-Yaniv, ``Selective Classification for Deep Neural
Networks,'' in \emph{NeurIPS}, 2017.

\bibitem{xin2021art}
J.~Xin et al., ``The Art of Abstention: Selective Prediction and Error
Regularization,'' \emph{JMLR}, 2021.

\bibitem{gal2016dropout}
Y.~Gal and Z.~Ghahramani, ``Dropout as a Bayesian Approximation,'' in
\emph{ICML}, 2016.

\bibitem{lakshminarayanan2017simple}
B.~Lakshminarayanan et al., ``Simple and Scalable Predictive Uncertainty
Estimation,'' in \emph{NeurIPS}, 2017.

\bibitem{wei2023inverse}
J.~Wei et al., ``Inverse Scaling Can Become U-Shaped,'' in \emph{EMNLP},
2023.

\bibitem{recht2019imagenet}
B.~Recht et al., ``Do ImageNet Classifiers Generalize to ImageNet?'' in
\emph{ICML}, 2019.

\bibitem{oquab2023dinov2}
M.~Oquab et al., ``DINOv2: Learning Robust Visual Features without
Supervision,'' \emph{TMLR}, 2023.

\bibitem{verma2021monusac}
R.~Verma et al., ``MoNuSAC 2020: A Multi-Organ Nuclei Segmentation and
Classification Challenge,'' in \emph{IEEE ISBI}, 2021.

\bibitem{vinyals2016matching}
O.~Vinyals et al., ``Matching Networks for One Shot Learning,'' in
\emph{NeurIPS}, 2016.

\bibitem{snell2017prototypical}
J.~Snell et al., ``Prototypical Networks for Few-shot Learning,'' in
\emph{NeurIPS}, 2017.

\bibitem{lei2024cross}
W.~Lei et al., ``Cross-Domain Few-Shot Segmentation via Iterative
Support-Query Correspondence Mining,'' in \emph{CVPR}, 2024.

\bibitem{shaban2017one}
A.~Shaban et al., ``One-Shot Learning for Semantic Segmentation,'' in
\emph{BMVC}, 2017.

\bibitem{nguyen2019feature}
K.~Nguyen and S.~Todorovic, ``Feature Weighting and Boosting for Few-Shot
Segmentation,'' in \emph{ICCV}, 2019.

\end{thebibliography}
}

\end{document}